\chapter{第7章：代数曲线 习题}
\difficultykey

\section{基础题 \easy}

\begin{problem}{射影化与无穷远点}
\easy
将仿射曲线
\[
E: y^2=x^3-x
\]
射影化到 $\mathbb{P}^2$ 中，并求它的全部无穷远点。
\end{problem}
\begin{solution}
令 $x=X/Z$、$y=Y/Z$，两边乘以 $Z^3$，得到射影方程
\[
Y^2Z=X^3-XZ^2.
\]
无穷远点满足 $Z=0$，于是方程化成 $X^3=0$，故 $X=0$。齐次坐标不能全为零，所以唯一无穷远点是
\[
[0:1:0].
\]
因此该曲线的射影闭包只有一个无穷远点，这正是椭圆曲线的单位元 $O$。
\end{solution}

\begin{problem}{光滑性与奇点}
\easy
判断下列曲线是否光滑；若不光滑，写出奇点：
\begin{itemize}
  \item[(i)] $y^2=x^3$；
  \item[(ii)] $y^2=x^2(x+1)$；
  \item[(iii)] $y^2=x^3-x$。
\end{itemize}
\end{problem}
\begin{solution}
设 $F(x,y)=0$ 为曲线方程。奇点条件是 $F=F_x=F_y=0$。
\begin{itemize}
  \item[(i)] $F=y^2-x^3$，故 $F_x=-3x^2$、$F_y=2y$。三式同时成立只可能在 $(0,0)$，所以它不光滑，奇点是 $(0,0)$。
  \item[(ii)] $F=y^2-x^3-x^2$，故 $F_x=-3x^2-2x=-x(3x+2)$、$F_y=2y$。由 $F_y=0$ 得 $y=0$，再由 $F=0$ 得 $x=0$ 或 $x=-1$。其中只有 $x=0$ 还满足 $F_x=0$，故奇点为 $(0,0)$。
  \item[(iii)] $F=y^2-x^3+x$，故 $F_x=-3x^2+1$、$F_y=2y$。若有奇点，则 $y=0$ 且 $x^3-x=0$，即 $x=0,\pm1$。但在这三点上 $F_x$ 分别为 $1,-2,-2$，都不为零，所以曲线光滑。
\end{itemize}
\end{solution}

\begin{problem}{基本不变量}
\easy
对短 Weierstrass 曲线
\[
E: y^2=x^3-x
\]
计算判别式 $\Delta$、$c_4$ 和 $j$-不变量。
\end{problem}
\begin{solution}
这里 $A=-1,B=0$。短 Weierstrass 形式的公式是
\[
\Delta=-16(4A^3+27B^2),\qquad c_4=-48A,\qquad j=\frac{c_4^3}{\Delta}.
\]
代入得
\[
\Delta=-16\cdot 4(-1)^3=64,
\qquad
c_4=48,
\qquad
j=\frac{48^3}{64}=1728.
\]
因此该曲线的判别式为 $64$，且 $j(E)=1728$。
\end{solution}

\begin{problem}{小素数处的点数}
\easy
手动计算曲线
\[
E: y^2=x^3-x
\]
在 $p=2,3,5,7$ 时的点数 $\#E(\mathbb{F}_p)$，并求相应的 $a_p=p+1-\#E(\mathbb{F}_p)$。
\end{problem}
\begin{solution}
逐个代入 $x\in\mathbb{F}_p$ 计数。

当 $p=2$ 时，$x^3-x\equiv 0$ 对两个 $x$ 都成立，因此仿射点为 $(0,0),(1,0)$，再加无穷远点，故
\[
\#E(\mathbb{F}_2)=3,
\qquad
a_2=2+1-3=0.
\]

当 $p=3$ 时，逐项计算可得仿射点为 $(0,0),(1,0),(2,0)$，故
\[
\#E(\mathbb{F}_3)=4,
\qquad
a_3=3+1-4=0.
\]

当 $p=5$ 时，平方剩余为 $0,1,4$。计算 $x^3-x$ 得
\[
0,0,1,4,0.
\]
于是仿射点数为 $1+1+2+2+1=7$，故
\[
\#E(\mathbb{F}_5)=8,
\qquad
a_5=5+1-8=-2.
\]

当 $p=7$ 时，平方剩余为 $0,1,2,4$。计算 $x^3-x$ 得
\[
0,0,6,3,4,1,0.
\]
只有 $x=0,1,4,5,6$ 有解，对应仿射点数 $1+1+2+2+1=7$，故
\[
\#E(\mathbb{F}_7)=8,
\qquad
a_7=7+1-8=0.
\]
因此
\[
(a_2,a_3,a_5,a_7)=(0,0,-2,0).
\]
\end{solution}

\begin{problem}{好约化与坏约化}
\easy
对同一条曲线 $E:y^2=x^3-x$，判断 $p=2,3,5,7$ 时是好约化还是坏约化。
\end{problem}
\begin{solution}
对短 Weierstrass 方程，坏约化恰发生在判别式被 $p$ 整除时。上一题中
\[
\Delta(E)=64=2^6.
\]
因此只有 $p=2$ 是坏约化，$p=3,5,7$ 都是好约化。

也可直接看模 $2$ 约化：
\[
y^2=x^3-x=x(x-1)(x+1)\equiv x(x+1)^2 \pmod 2,
\]
出现重根，所以 $p=2$ 的约化不光滑。其余三个素数不整除 $\Delta$，故约化光滑。
\end{solution}

\begin{problem}{导子十一曲线的不变量}
\easy
设
\[
E_{11}: y^2+y=x^3-x^2-10x-20.
\]
求它的 $b_2,b_4,b_6,b_8$、$c_4$ 与判别式 $\Delta$。
\end{problem}
\begin{solution}
把方程写成一般 Weierstrass 形式
\[
y^2+a_1xy+a_3y=x^3+a_2x^2+a_4x+a_6,
\]
可读出
\[
a_1=0,\ a_2=-1,\ a_3=1,\ a_4=-10,\ a_6=-20.
\]
因此
\[
b_2=a_1^2+4a_2=-4,
\qquad
b_4=2a_4+a_1a_3=-20,
\qquad
b_6=a_3^2+4a_6=-79,
\]
\[
b_8=a_1^2a_6+4a_2a_6-a_1a_3a_4+a_2a_3^2-a_4^2=-21.
\]
再用公式
\[
c_4=b_2^2-24b_4,
\qquad
\Delta=-b_2^2b_8-8b_4^3-27b_6^2+9b_2b_4b_6,
\]
得到
\[
c_4=16+480=496,
\qquad
\Delta=-161051=-11^5.
\]
所以这条曲线唯一的坏素数是 $11$。
\end{solution}

\begin{problem}{导子十一曲线的小素数点数}
\easy
计算 $E_{11}$ 在 $p=2,3,5,7$ 时的点数和 $a_p$。
\end{problem}
\begin{solution}
直接代入计算。

当 $p=2$ 时，方程化为
\[
y^2+y=x^3+x^2.
\]
枚举 $x,y\in\mathbb{F}_2$ 可得 $4$ 个仿射点，再加无穷远点，所以
\[
\#E_{11}(\mathbb{F}_2)=5,
\qquad a_2=2+1-5=-2.
\]

当 $p=3$ 时，逐项检查得到同样有 $4$ 个仿射点，因此
\[
\#E_{11}(\mathbb{F}_3)=5,
\qquad a_3=3+1-5=-1.
\]

当 $p=5$ 时，有 $4$ 个仿射点，所以
\[
\#E_{11}(\mathbb{F}_5)=5,
\qquad a_5=5+1-5=1.
\]

当 $p=7$ 时，有 $9$ 个仿射点，所以
\[
\#E_{11}(\mathbb{F}_7)=10,
\qquad a_7=7+1-10=-2.
\]
因此
\[
(a_2,a_3,a_5,a_7)=(-2,-1,1,-2).
\]
\end{solution}

\begin{problem}{哈斯界的检验}
\easy
用上题结果检验 $E_{11}$ 在 $p=2,3,5,7,13$ 时满足 Hasse 界
\[
|a_p|\le 2\sqrt{p}.
\]
其中可用已知值 $a_{13}=4$。
\end{problem}
\begin{solution}
逐个比较即可：
\[
|a_2|=2\le 2\sqrt2,
\qquad
|a_3|=1\le 2\sqrt3,
\qquad
|a_5|=1\le 2\sqrt5,
\]
\[
|a_7|=2\le 2\sqrt7,
\qquad
|a_{13}|=4\le 2\sqrt{13}.
\]
全部成立。这也等价于
\[
p+1-2\sqrt p\le \#E_{11}(\mathbb{F}_p)\le p+1+2\sqrt p.
\]
\end{solution}

\begin{problem}{Hecke 对应的代数描述}
\easy
设 $p\nmid N$。说明模曲线 $X_0(N)$ 上的 Hecke 对应 $T_p$ 可以用 $X_0(Np)$ 到 $X_0(N)$ 的两条退化映射来描述。
\end{problem}
\begin{solution}
$X_0(Np)$ 的点可看成三元组 $(E,C_N,C_p)$，其中 $C_N$ 是 $N$ 阶循环子群，$C_p$ 是与之相容的 $p$ 阶循环子群。于是有两条自然映射
\[
\alpha(E,C_N,C_p)=(E,C_N),
\]
和
\[
\beta(E,C_N,C_p)=\bigl(E/C_p,(C_N+C_p)/C_p\bigr).
\]
Hecke 对应正是由这两条映射给出的有限对应
\[
X_0(N) \xleftarrow{\ \alpha\ } X_0(Np) \xrightarrow{\ \beta\ } X_0(N).
\]
在除子类群或 Jacobian 上，它诱导出
\[
T_p=\beta_*\alpha^*=
\alpha_*\beta^*.
\]
这就是 Hecke 算子的代数定义。
\end{solution}

\section{进阶题 \medium}

\begin{problem}{十一处的奇点与结点}
\medium
证明 $E_{11}$ 在 $p=11$ 处是坏约化，并找出模 $11$ 约化后的奇点；再说明这个奇点是结点而不是尖点。
\end{problem}
\begin{solution}
由上题 $\Delta(E_{11})=-11^5$，可知 $11$ 是坏素数。把方程模 $11$ 化简得
\[
y^2+y=x^3-x^2+x+2.
\]
设
\[
F(x,y)=y^2+y-x^3+x^2-x-2.
\]
则
\[
F_x=-3x^2+2x-1,
\qquad
F_y=2y+1.
\]
在 $\mathbb{F}_{11}$ 中，由 $F_y=0$ 得 $y=5$。再代入 $F_x=0$ 得
\[
3x^2-2x+1=0,
\]
检验可得解 $x=5$，并且 $F(5,5)=0$。所以唯一奇点是 $(5,5)$。

令 $x=u+5$、$y=v+5$，则局部方程化为
\[
v^2=u^2(u+3).
\]
最低次齐次项是
\[
v^2-3u^2=(v-5u)(v+5u)
\]
，因为 $5^2=25\equiv 3\pmod{11}$。它分解成两条不同直线，所以该奇点是结点。因此 $E_{11}$ 在 $11$ 处是乘法坏约化。
\end{solution}

\begin{problem}{塔玛伽瓦数的计算}
\medium
已知 $E_{11}$ 的 N\'eron 最小模型在 $p=11$ 的特殊纤维是 Kodaira 型 $I_5$。求 Tamagawa 数 $c_{11}$。
\end{problem}
\begin{solution}
对乘法型约化，Kodaira 型 $I_n$ 的特殊纤维由 $n$ 个不可约分量构成一个环。分量群的阶就是相应的 Tamagawa 数。因此当特殊纤维型为 $I_5$ 时，
\[
c_{11}=5.
\]
这与 $\ord_{11}(\Delta)=5$ 相符：最小判别式的 $11$-进赋值正好给出乘法型中的指标 $n$。
\end{solution}

\begin{problem}{非最小模型的缩放}
\medium
考虑曲线
\[
E': y^2=x^3-16x.
\]
通过变量代换把它化成一个更小的整系数模型，并比较两者的判别式。
\end{problem}
\begin{solution}
令
\[
x=4u,\qquad y=8v.
\]
代入得
\[
64v^2=64u^3-64u,
\]
即
\[
v^2=u^3-u.
\]
所以 $E'$ 与
\[
E: v^2=u^3-u
\]
同构，而后者的系数更小，是自然的最小模型。

对短 Weierstrass 方程 $y^2=x^3+Ax+B$，判别式为 $-16(4A^3+27B^2)$。于是
\[
\Delta(E')=-16\bigl(4(-16)^3\bigr)=2^{16},
\qquad
\Delta(E)=64=2^6.
\]
二者相差 $2^{10}=2^{12}/2^2$，这正反映了缩放 $x=u\,2^2$、$y=v\,2^3$ 会把判别式乘上 $2^{12}$ 的逆向因子。故 $E'$ 不是最小模型，而 $E$ 更接近最小。
\end{solution}

\begin{problem}{零迹现象}
\medium
设 $E:y^2=x^3-x$。证明当奇素数 $p\equiv 3\pmod 4$ 时，$a_p(E)=0$。
\end{problem}
\begin{solution}
当 $p\equiv 3\pmod 4$ 时，$-1$ 不是平方剩余。对每个 $x\in\mathbb{F}_p$，方程
\[
y^2=x^3-x=x(x-1)(x+1)
\]
的右端记为 $r_x$。

若 $r_x\ne 0$，则有解当且仅当 $r_x$ 是平方剩余；此时恰有两个解 $y=\pm\sqrt{r_x}$。注意把 $x$ 换成 $-x$ 时，右端变成
\[
(-x)^3-(-x)=-(x^3-x)=-r_x.
\]
由于 $-1$ 不是平方剩余，$r_x$ 与 $-r_x$ 中恰有一个是平方剩余。因此 $x$ 与 $-x$ 成对贡献总共 $2$ 个仿射点。

再看不动点 $x=0,\pm1$，这三点都使右端为 $0$，各贡献一个仿射点。于是总仿射点数为
\[
(p-3)+3=p,
\]
再加无穷远点得
\[
\#E(\mathbb{F}_p)=p+1.
\]
因此
\[
a_p(E)=p+1-\#E(\mathbb{F}_p)=0.
\]
\end{solution}

\begin{problem}{哈斯区间的应用}
\medium
利用 Hasse 界说明：若已知某条椭圆曲线在 $\mathbb{F}_{13}$ 上有 $10$ 个点，则其 $a_{13}$ 是多少；并说明这与 Hasse 界相容。
\end{problem}
\begin{solution}
由定义
\[
a_{13}=13+1-10=4.
\]
而 Hasse 界给出
\[
|a_{13}|\le 2\sqrt{13}\approx 7.21.
\]
显然 $|4|<7.21$，所以完全相容。换句话说，$10$ 个点落在允许区间
\[
13+1-2\sqrt{13}\le \#E(\mathbb{F}_{13})\le 13+1+2\sqrt{13}
\]
之内。
\end{solution}

\begin{problem}{有理二挠点}
\medium
求曲线 $E:y^2=x^3-x$ 的全部有理二挠点，并说明它们构成什么群。
\end{problem}
\begin{solution}
在一般 Weierstrass 方程中，$2$-挠点满足 $y=0$。因此只需求
\[
x^3-x=x(x-1)(x+1)=0.
\]
于是三个有限点是
\[
(0,0),\ (1,0),\ (-1,0).
\]
再加无穷远点 $O$，得到
\[
E[2](\Q)=\{O,(0,0),(1,0),(-1,0)\}.
\]
这个集合有四个元素，而且每个非零元都满足两倍为零，所以
\[
E[2](\Q)\cong (\Z/2\Z)^2.
\]
\end{solution}

\begin{problem}{导体的确定}
\medium
用约化信息说明 $E_{11}$ 的导体是 $N_{E_{11}}=11$。
\end{problem}
\begin{solution}
导体只在坏素数处有贡献。由判别式
\[
\Delta(E_{11})=-11^5
\]
可知只有 $11$ 是坏素数。上一题已经说明在 $11$ 处是乘法坏约化。对乘法坏约化，导体指数等于 $1$，所以
\[
N_{E_{11}}=11^1=11.
\]
也就是说，$E_{11}$ 的导体恰好等于唯一坏素数本身。
\end{solution}

\begin{problem}{特殊纤维与分量群}
\medium
解释为什么 $E_{11}$ 在 $11$ 处的特殊纤维不是不可约的，并由此说明分量群是五阶循环群。
\end{problem}
\begin{solution}
因为 $E_{11}$ 在 $11$ 处是乘法坏约化，N\'eron 最小模型的特殊纤维不再是一个光滑椭圆曲线，而是若干条射影直线首尾相接形成的环。已知其 Kodaira 型为 $I_5$，所以共有 $5$ 个不可约分量。

把特殊纤维的单位分量记作 $\mathscr{E}^0$。分量群定义为全部分量模掉单位分量，故其阶等于分量个数；于是
\[
\Phi_{11}\cong \Z/5\Z.
\]
因此 Tamagawa 数 $c_{11}=\#\Phi_{11}(\mathbb{F}_{11})=5$。
\end{solution}

\begin{problem}{对应的次数}
\medium
设 $p\nmid N$。说明 Hecke 对应
\[
X_0(N) \xleftarrow{\ \alpha\ } X_0(Np) \xrightarrow{\ \beta\ } X_0(N)
\]
中，两条映射都具有次数 $p+1$。
\end{problem}
\begin{solution}
固定 $X_0(N)$ 上一点 $(E,C_N)$。要描述 $\alpha^{-1}(E,C_N)$，只需在 $E$ 上选取一个与 $C_N$ 相容的 $p$ 阶循环子群 $C_p$。椭圆曲线上这样的子群恰有 $p+1$ 个，因此
\[
\deg \alpha=p+1.
\]

对 $\beta$ 也是一样：固定 $(E',C_N')$，要找它的原像，相当于寻找所有满足
\[
E/C_p\cong E'
\]
的 $p$ 阶子群 $C_p$。这同样对应于 $E'[p]$ 中的一维子空间，也有 $p+1$ 个，所以
\[
\deg \beta=p+1.
\]
因此 Hecke 对应在两侧都是次数 $p+1$ 的有限对应。
\end{solution}

\section{挑战题 \hard}

\begin{problem}{数值乘积的证据}
\hard
对 $E_{11}$ 定义
\[
P(X)=\prod_{p\le X,\ p\ne 11}\frac{p}{\#E_{11}(\mathbb{F}_p)}.
\]
利用已知点数计算 $P(5)$、$P(13)$、$P(29)$，并解释为什么这给出 BSD 猜想的秩零数值证据。
\end{problem}
\begin{solution}
由前面计算，
\[
\#E_{11}(\mathbb{F}_2)=5,
\ \#E_{11}(\mathbb{F}_3)=5,
\ \#E_{11}(\mathbb{F}_5)=5,
\ \#E_{11}(\mathbb{F}_7)=10,
\ \#E_{11}(\mathbb{F}_{13})=10.
\]
再利用已知数据
\[
\#E_{11}(\mathbb{F}_{17})=20,
\quad
\#E_{11}(\mathbb{F}_{19})=20,
\quad
\#E_{11}(\mathbb{F}_{23})=25,
\quad
\#E_{11}(\mathbb{F}_{29})=30,
\]
得到
\[
P(5)=\frac25\cdot\frac35\cdot\frac55=0.24,
\]
\[
P(13)=P(5)\cdot\frac7{10}\cdot\frac{13}{10}=0.2184,
\]
\[
P(29)=P(13)\cdot\frac{17}{20}\cdot\frac{19}{20}\cdot\frac{23}{25}\cdot\frac{29}{30}\approx 0.15684.
\]
这些值没有表现出明显增长，更像是围绕常数波动。BSD 的秩零情形预言相关 Euler 型乘积应接近常数量级，因此这是支持 $\operatorname{rank}E_{11}(\Q)=0$ 的数值证据。
\end{solution}

\begin{problem}{模曲线的有理点}
\hard
利用第 6 章已经知道的同构 $X_0(11)\cong E_{11}$ 以及 $E_{11}(\Q)\cong \Z/5\Z$，说明 $X_0(11)(\Q)$ 恰有五个有理点。
\end{problem}
\begin{solution}
既然第 6 章已经把 $X_0(11)$ 与椭圆曲线 $E_{11}$ 识别起来，那么两者的有理点集合等同：
\[
X_0(11)(\Q)\cong E_{11}(\Q).
\]
又已知
\[
E_{11}(\Q)\cong \Z/5\Z,
\]
所以它的有理点总数正好是 $5$。因此
\[
\#X_0(11)(\Q)=5.
\]
这五个点中包含两个 cusp，其余三个点对应非 cusp 的有理 $11$-同源数据。
\end{solution}

\begin{problem}{可见的五个有理点}
\hard
验证 $E_{11}$ 上下列点都在曲线上：
\[
O,\ (5,5),\ (5,-6),\ (16,60),\ (16,-61),
\]
并说明这五个点构成一个五阶循环群。
\end{problem}
\begin{solution}
把 $x=5$ 代入方程右端得
\[
5^3-5^2-10\cdot 5-20=30,
\]
而
\[
5^2+5=30,
\qquad
(-6)^2+(-6)=30,
\]
所以 $(5,5)$ 和 $(5,-6)$ 都在曲线上。再把 $x=16$ 代入得
\[
16^3-16^2-10\cdot16-20=3780,
\]
而
\[
60^2+60=3780,
\qquad
(-61)^2+(-61)=3780,
\]
故 $(16,60)$、$(16,-61)$ 也在曲线上。

第 6 章已经说明 $E_{11}(\Q)\cong \Z/5\Z$。既然我们已经显式找到了 $O$ 之外的四个有理点，而整个群恰有五个元素，所以这五个点正好就是全部有理点，并自动构成一个五阶循环群。
\end{solution}

\begin{problem}{挠点与模二模三约化}
\hard
利用约化映射说明 $E_{11}(\Q)_{\mathrm{tors}}$ 的阶必须整除 $5$，再与上一题结合推出
\[
E_{11}(\Q)_{\mathrm{tors}}\cong \Z/5\Z.
\]
\end{problem}
\begin{solution}
当 $p$ 是好素数时，阶与 $p$ 互素的有理挠点可嵌入到 $E_{11}(\mathbb{F}_p)$ 中。因此挠子群的阶必须同时整除若干好约化点群的阶。

这里
\[
\#E_{11}(\mathbb{F}_2)=5,
\qquad
\#E_{11}(\mathbb{F}_3)=5,
\qquad
\#E_{11}(\mathbb{F}_5)=5.
\]
于是任何与 $2,3,5$ 互素的挠点阶都必须整除 $5$；实际上整个挠子群阶也只能整除这些群阶的最大公因数，故必整除 $5$。上一题已经找到了一个五阶有理子群，所以只能有
\[
E_{11}(\Q)_{\mathrm{tors}}\cong \Z/5\Z.
\]
因为第 6 章又说明 $E_{11}(\Q)$ 没有自由部分，这其实就是整个 Mordell--Weil 群。
\end{solution}

\begin{problem}{分裂乘法约化}
\hard
证明 $E_{11}$ 在 $11$ 处不只是乘法约化，而且是分裂乘法约化。
\end{problem}
\begin{solution}
在第一个进阶题中，我们把模 $11$ 的奇点平移到原点后得到局部方程
\[
v^2=u^2(u+3).
\]
其最低次齐次部分分解为
\[
v^2-3u^2=(v-5u)(v+5u)
\]
，因为 $5^2\equiv 3\pmod{11}$。这说明结点的两条切线都已经定义在 $\mathbb{F}_{11}$ 上，而不是要到二次扩张里才出现。

结点若在底域上分裂成两条不同切线，就对应分裂乘法约化；若只在二次扩张中分裂，则是非分裂乘法约化。因此 $E_{11}$ 在 $11$ 处是分裂乘法约化。
\end{solution}

\begin{problem}{Hecke 对应对除子类的作用}
\hard
设 $p\nmid N$，$D$ 是 $X_0(N)$ 上的一个次数为零的有理除子。说明 Hecke 对应把 $D$ 送到另一个次数为零的有理除子类。
\end{problem}
\begin{solution}
Hecke 对应由有限映射
\[
X_0(N) \xleftarrow{\ \alpha\ } X_0(Np) \xrightarrow{\ \beta\ } X_0(N)
\]
给出，因此对除子有操作
\[
T_p(D)=\beta_*\alpha^*(D).
\]
先看次数。因为拉回会把每个点替换成 $\deg\alpha=p+1$ 个点的和，所以
\[
\deg(\alpha^*D)=(p+1)\deg D=0.
\]
再推前时，次数保持不变，因此
\[
\deg(T_pD)=0.
\]

若 $D$ 在 $\Q$ 上有理，则 $\alpha,\beta$ 本身也在 $\Q$ 上定义，所以 $\alpha^*D$ 和 $\beta_*\alpha^*D$ 仍然在 $\Q$ 上有理。于是 $T_p$ 保持 $\mathrm{Pic}^0(X_0(N))$，从而在 Jacobian 上给出良定义的自同态。
\end{solution}

\begin{problem}{模曲线的小素数点数}
\hard
利用 $X_0(11)\cong E_{11}$ 以及第 7 章前面算出的 $a_2,a_3,a_5,a_7,a_{13}$，计算
\[
\#X_0(11)(\mathbb{F}_p)
\]
在 $p=2,3,5,7,13$ 时的值。
\end{problem}
\begin{solution}
因为 $X_0(11)$ 在 $\mathbb{Q}$ 上与椭圆曲线 $E_{11}$ 同构，所以对每个好素数 $p\ne 11$ 都有
\[
\#X_0(11)(\mathbb{F}_p)=\#E_{11}(\mathbb{F}_p)=p+1-a_p.
\]
于是
\[
\#X_0(11)(\mathbb{F}_2)=2+1-(-2)=5,
\qquad
\#X_0(11)(\mathbb{F}_3)=3+1-(-1)=5,
\]
\[
\#X_0(11)(\mathbb{F}_5)=5+1-1=5,
\qquad
\#X_0(11)(\mathbb{F}_7)=7+1-(-2)=10,
\]
\[
\#X_0(11)(\mathbb{F}_{13})=13+1-4=10.
\]
因此这些素数处的点数分别是
\[
5,\ 5,\ 5,\ 10,\ 10.
\]
这个计算说明模曲线的有限域点数与相应 Hecke 系数完全同步，这正是后续 Eichler--Shimura 关系会系统解释的现象。
\end{solution}

\begin{problem}{最小判别式与最小模型}
\hard
设一条整系数 Weierstrass 方程经过变量代换
\[
x=u^2x',\qquad y=u^3y'
\]
化成另一条整系数方程。说明判别式会缩放为 $\Delta'=u^{-12}\Delta$，并据此解释为什么最小模型追求的是使 $|\Delta|$ 尽量小。
\end{problem}
\begin{solution}
在 Weierstrass 方程的标准变量代换中，$x$ 具有权 $2$，$y$ 具有权 $3$，因此判别式作为权 $12$ 的量，会按
\[
\Delta\longmapsto u^{-12}\Delta
\]
缩放。这一点也可由短 Weierstrass 形式直接检查：若
\[
y^2=x^3+Ax+B,
\]
则代换后系数变成 $A'=u^{-4}A$、$B'=u^{-6}B$，于是
\[
\Delta'=-16(4A'^3+27B'^2)=u^{-12}\Delta.
\]
因此如果还能继续把方程除去一个公共的 $u$，判别式的绝对值就会明显减小。所谓最小模型，就是在所有整系数模型中把这种可缩放性用尽，使得局部判别式赋值不再下降；全局上就得到绝对值最小的判别式。
\end{solution}

\begin{problem}{从模曲线到椭圆曲线模型}
\hard
综合说明为什么 $X_0(11)$ 能在 $\mathbb{Q}$ 上写成椭圆曲线 $E_{11}$，并解释导体 $11$、坏约化点 $11$ 与 Hecke 对应三者之间的关系。
\end{problem}
\begin{solution}
第 5 章的维数公式给出
\[
g_0(11)=\dim S_2(\Gamma_0(11))=1.
\]
因此 $X_0(11)$ 是一条亏格 $1$ 的模曲线。第 6 章进一步把它的 Jacobian 与它自身识别，于是得到一条定义在 $\mathbb{Q}$ 上的椭圆曲线模型；常用最小模型正是
\[
E_{11}: y^2+y=x^3-x^2-10x-20.
\]

由判别式
\[
\Delta(E_{11})=-11^5
\]
知只有 $11$ 是坏素数，所以导体也只能由 $11$ 贡献；又因为模 $11$ 的奇点是结点，所以是乘法坏约化，导体指数为 $1$，从而导体恰为 $11$。

另一方面，对每个 $p\nmid 11$，Hecke 对应 $T_p$ 在模曲线上有明确的代数定义。它们控制了 Jacobian 上的自同态结构，也正是后来 Eichler--Shimura 关系的几何起点。于是这条代数曲线既记录了 level $11$ 的模意义，又带有与 Hecke 算子相容的算术结构。
\end{solution}
